\section{Actual affine-localization stalk and residue maps}

\begin{theorem}[Original affine localization gives native stalk isomorphisms]
\label{thm:fc-affine-localization-native-stalks}
\leanok
\lean{EnvelopeProperty.AffineLocalizationStalks.stalkMap_isIso}
For an actual map $g:Y\to X$ of actual affine schemes whose literal original
whole section map is a localization, every same original native stalk map
of $g$ is an isomorphism.
\uses{thm:fc-spec-localization-native-stalk-residue}
\end{theorem}
\begin{proof}
The native whole affine-spectrum identifications carry the original $g$ to
the spectrum of its literal original section map.  The spectrum-localization
stalk theorem and the actual whole naturality square give isomorphisms of
the original composite stalk maps.  Cancel the original affine-spectrum
stalk isomorphisms.  The only generic input is that the actual section ring
map is a localization; no stalk comparison or point correspondence is given.
\end{proof}

\begin{theorem}[Original affine-localization residue maps]
\label{thm:fc-affine-localization-native-residue-fields}
\leanok
\lean{EnvelopeProperty.AffineLocalizationStalks.residueFieldMap_isIso}
Every same original native residue-field map of that original $g$ is an
isomorphism, preserving its original point and coefficient map.
\uses{thm:fc-affine-localization-native-stalks}
\end{theorem}
\begin{proof}
Apply the native residue-field construction to the actual original stalk
isomorphism.  The concrete localized-inverse consumer must derive the
whole original section-localization property internally; this generic
lemma does not supply a desired finite-fibre degree formula.
\end{proof}
